\chapter{Réalisation technique}
\label{ch:real}

\section{Composition racine et pipeline}

Le fichier \texttt{Program.cs} de l'application Web enregistre Blazor (composants interactifs serveur), le \texttt{DbContext} SQLite, et les services métier via une méthode d'extension :

\begin{lstlisting}[language=CSharp, caption={Extrait typique d'enregistrement des services}]
builder.Services.AddDbContext<AppDbContext>(options =>
    options.UseSqlite(builder.Configuration.GetConnectionString("DefaultConnection")));
builder.Services.AddPortfolioServices(builder.Configuration);
\end{lstlisting}

Après construction du pipeline (fichiers statiques, antiforgery, Razor Components), une portée \texttt{scoped} exécute les migrations et le \emph{seed}.

\section{Services métier}

\subsection{Portefeuille (\texttt{PortfolioService})}

Calcule l'instantané \texttt{PortfolioSnapshot} à partir des actifs, des transactions et du budget. Le \textbf{coût moyen pondéré} (CMP) des titres encore détenus est obtenu par une passe chronologique sur les opérations (\texttt{PositionCostHelper}) : à l'achat, le CMP est recalculé ; à la vente partielle, le CMP unitaire des titres restants est conservé (logique courante pour une valorisation de portefeuille pédagogique).

\subsection{Transactions (\texttt{TransactionService})}

Vérifie la liquidité pour les achats et la quantité détenue pour les ventes, met à jour \texttt{PortfolioBudget.AvailableCash}, persiste la transaction, puis enregistre un point d'historique de valeur totale via \texttt{IPortfolioService}.

\subsection{Actifs (\texttt{AssetService})}

CRUD, filtrage LINQ (recherche insensible à la casse sur symbole/nom, filtre par type), pagination serveur (\texttt{PagedResult<T>}).

\subsection{Prix marché (\texttt{PriceDataService})}

Selon \texttt{PriceData:Mode} (\texttt{Simulated}, \texttt{Api}, \texttt{None}), met à jour ou non les \texttt{CurrentPrice}. En mode API, utilisation de Finnhub (actions, clé API) et CoinGecko (cryptos mappées), avec mise en cache mémoire pour limiter les appels.

\subsection{Alertes (\texttt{PortfolioAlertService})}

Agrège instantané et analytique pour produire une liste de \texttt{PortfolioAlert} selon les seuils \texttt{appsettings.json}.

\section{Validation}

Les entités \texttt{Asset}, \texttt{Transaction}, \texttt{PortfolioBudget}, etc. portent des attributs \texttt{[Required]}, \texttt{[Range]}, \texttt{[StringLength]}. Les formulaires Blazor utilisent \texttt{EditForm}, \texttt{DataAnnotationsValidator} et \texttt{ValidationSummary}.

\section{Asynchronisme}

Toutes les opérations d'accès base dans les services sont asynchrones (\texttt{Task}, \texttt{async}/\texttt{await}), ce qui évite de bloquer les threads du pool lors des E/S réseau et base de données.
